\section{La voz como biomarcador. Fundamentos}\label{sec:voz_biomarcador}

La voz humana es el resultado de un proceso fisiológico y biomecánico altamente complejo que exige la coordinación síncrona de múltiples subsistemas: el aparato respiratorio (fuente aerodinámica), el sistema fonatorio (vibración cordal), el sistema de resonancia y el aparato articulatorio, todo ello bajo un estricto control neuromotor y cognitivo \cite{romero2018}. Cualquier alteración sistémica en estos mecanismos, especialmente aquellas derivadas de patologías que comprometen el sistema nervioso central, induce alteraciones acústicas sutiles pero matemáticamente cuantificables. Por consiguiente, la señal de voz adquiere la categoría de biomarcador digital para el estudio y la detección temprana de enfermedades neurodegenerativas.

A lo largo de las últimas décadas, la literatura clínica ha evidenciado empíricamente que estas disfunciones vocales se manifiestan en fases preclínicas, precediendo a los síntomas motores o cognitivos evidentes. En la Enfermedad de Parkinson, la rigidez muscular y la bradicinesia afectan directamente a la laringe, generando hipofonía (reducción del volumen), monotonía prosódica e inestabilidad vibratoria, cuantificable a través del incremento en la varianza de la frecuencia (\textit{jitter}) y de la amplitud (\textit{shimmer}) \cite{moro-velazquezNewToolsDifferential2021}. En la Enfermedad de Alzheimer, el deterioro neurocognitivo se proyecta en el discurso mediante un incremento en la frecuencia y duración de las pausas silenciosas, una reducción drástica de la fluidez verbal, aplanamiento de la curva melódica y alteraciones en la velocidad de articulación \cite{garcia-gutierrezUnveilingSoundCognitive2024}. Por su parte, en la variante bulbar de la Esclerosis Lateral Amiotrófica (ELA), la degeneración neuromuscular desencadena disartria, hipernasalidad y fatiga respiratoria prematura \cite{gutz2019}.

La principal ventaja de la adopción de la voz como biomarcador radica en su naturaleza intrínsecamente no invasiva y en su bajo coste de adquisición. Las bioseñales acústicas pueden ser capturadas con micrófonos o dispositivos móviles convencionales, democratizando el acceso a las pruebas de cribado sin requerir infraestructura clínica especializada. Asimismo, la digitalización de la señal permite una automatización total del análisis computacional mediante algoritmos de procesamiento digital de señales (DSP) e Inteligencia Artificial, habilitando una monitorización epidemiológica masiva, continua y remota.

Desde el punto de vista metodológico, la extracción de estos biomarcadores requiere el diseño de protocolos que incluyan diversas tareas fonatorias, orientadas a aislar subsistemas neuromotores específicos:

\begin{itemize}
    \item \textbf{Fonación sostenida de vocales} (ej. la vocal /a/ prolongada): Elimina la variable articulatoria y lingüística, permitiendo evaluar de forma pura la estabilidad aerodinámica, la vibración mecánica de las cuerdas vocales y la armonicidad de la señal.
    \item \textbf{Lectura o repetición de frases estandarizadas:} Minimiza la carga cognitiva para centrar el análisis en la precisión articulatoria, la transición fonética, el control de la microprosodia y la coordinación velofaríngea.
    \item \textbf{Habla espontánea (monólogos libres):} Impone la mayor carga cognitiva y lingüística, facilitando la extracción de métricas complejas asociadas a la planificación del discurso, la fluidez, la sintaxis y la riqueza semántica.
\end{itemize}

A nivel técnico, la parametrización de las señales acústicas se automatiza empleando herramientas e interfaces de software ampliamente validadas por la comunidad científica, tales como \textit{openSMILE} \cite{eyben2010}, \textit{librosa} \cite{mcfee2015} o \textit{Praat} \cite{boersma2001}. Estas librerías permiten la extracción sistemática de centenares de características acústicas (temporales, espectrales y cepstrales) que, posteriormente, actúan como vectores de entrada (\textit{feature vectors}) para el entrenamiento de algoritmos de clasificación.

En síntesis, la voz constituye una fuente biométrica rica, robusta y altamente sensible a las alteraciones neuromotoras. Su integración con técnicas de Inteligencia Artificial representa un cambio de paradigma hacia el desarrollo de sistemas de apoyo a la decisión clínica (CDSS) objetivos y escalables.